\documentclass[12pt, a4paper]{report}
\usepackage[utf8]{inputenc}
\usepackage[T1]{fontenc}
\usepackage[french]{babel}
\usepackage{graphicx}
\usepackage{hyperref}
\usepackage{geometry}
\usepackage{fancyhdr}
\usepackage{listings}
\usepackage{xcolor}
\usepackage{lipsum} % Module indispensable qui genere du vrai-faux texte pour etendre la taille du document

% Configuration des marges
\geometry{top=2.5cm, bottom=2.5cm, left=3cm, right=2.5cm}

% Configuration des couleurs pour le code source
\definecolor{codegreen}{rgb}{0,0.6,0}
\definecolor{codegray}{rgb}{0.5,0.5,0.5}
\definecolor{codepurple}{rgb}{0.58,0,0.82}
\definecolor{backcolour}{rgb}{0.95,0.95,0.92}

\lstset{
    backgroundcolor=\color{backcolour},   
    commentstyle=\color{codegreen},
    keywordstyle=\color{magenta},
    numberstyle=\tiny\color{codegray},
    stringstyle=\color{codepurple},
    basicstyle=\ttfamily\footnotesize,
    breakatwhitespace=false,         
    breaklines=true,                 
    captionpos=b,                    
    keepspaces=true,                 
    numbers=left,                    
    numbersep=5pt,                  
    showspaces=false,                
    showstringspaces=false,
    showtabs=false,                  
    tabsize=2
}

\begin{document}

\begin{titlepage}
    \centering
    \vspace*{2cm}
    
    {\scshape\Large Rapport de Projet de Fin d'Études \par}
    \vspace{1.5cm}
    {\huge\bfseries Conception et Développement Full-Stack d'une Plateforme IA Data Analyst\par}
    \vspace{2cm}
    {\Large\itshape Présenté par : Akram Mouhlal\par}
    \vfill
    {\large Année Universitaire 2025 - 2026 \par}
\end{titlepage}

\tableofcontents
\listoffigures
\newpage

% --- INTRODUCTION ---
\chapter*{Introduction Générale}
\addcontentsline{toc}{chapter}{Introduction Générale}
Dans un monde où les données connaissent une croissance exponentielle, le domaine de la Data Science s'impose comme un levier stratégique majeur. Cependant, la complexité des algorithmes d'intelligence artificielle et le pré-traitement des données représentent des freins pour les utilisateurs non-techniques. 
L'objectif de ce projet est de concevoir une application web "Full-Stack" ergonomique, intelligente et automatisée, capable de traiter, nettoyer et analyser des fichiers volumineux (CSV) afin de générer des modèles prédictifs et descriptifs.

\lipsum[1-8] % Texte genere pour gonfler le volume du rapport - Remplacez par vos recherches

% --- CHAPITRE 1 ---
\chapter{Contexte du Projet et Cahier des Charges}
\section{Mise en contexte}
L'essor du Machine Learning nécessite des plateformes intuitives. Ce projet vise à réduire la friction mathématique entre un fichier brut et une analyse décisionnelle compréhensible.
\lipsum[9-18]

\section{Objectifs du projet}
Notre produit s'articule autour de trois axes de développement majeurs :
\begin{itemize}
    \item \textbf{Exploration Automatique (\textit{EDA}) :} Automatiser la découverte des données via la bibliothèque \texttt{pandas}.
    \item \textbf{Modélisation IA :} Mettre à disposition des algorithmes d'apprentissage complexes (K-Means, Régression Linéaire) contrôlables via des commandes React fluides.
    \item \textbf{Reporting Dynamique :} Exporter et formater les tableaux de bord interactifs en format PDF de manière native sans perte de fluidité.
\end{itemize}
\lipsum[19-28]

% --- CHAPITRE 2 ---
\chapter{État de l'Art et Choix Technologiques}
\section{L'Architecture Full-Stack et le Pattern API}
L'architecture choisie repose sur une séparation stricte entre le Frontend (L'interface visuelle) et le Backend (Le moteur de données) selon un modèle REST.
\lipsum[29-38]

\subsection{Frontend : React, Vite \& TailwindCSS}
L'interface a été conçue de manière moderne et haut de gamme, utilisant le style "Dark Mode" et le "Glassmorphism" (composants translucides floutés). \textit{React} (propulsé par Vite pour une vitesse fulgurante) a été sélectionné pour sa gestion optimisée et réactive du DOM.
\lipsum[39-44]

\subsection{Backend : Django REST Framework (Python)}
Le traitement intensif des données requiert Python. Le framework \textit{Django} s'est naturellement imposé. L'association au DRF permet d'exposer les exécutions asynchrones des modèles sous forme de routes API (JSON).
\lipsum[45-52]

\subsection{Boîte à outil d'Intelligence Artificielle}
Le moteur analytique a exploité principalement \textit{Scikit-Learn} et \textit{Numpy}.
\lipsum[53-58]

% --- CHAPITRE 3 ---
\chapter{Analyse et Modélisation}
\section{Modélisation de la Base de Données}
Le modèle relationnel (ORM Django) s'opère sur la table clé "Dataset" reliant chaque utilisateur au fichier CSV qu'il traite. La persistance inclut un résumé JSON complet de l'extraction des colonnes.
\lipsum[59-69]

\section{Architecture Séquentielle Systémique}
Parcours typique et cycle de vie d'une analyse :
\begin{enumerate}
    \item Le client (React) soumet via Axios un fichier au format MultiPart Form Data.
    \item Le backend Python stocke de manière sécurisée l'archive.
    \item Un parseur d'intégrité Pandas s'active, compte les valeurs, supprime les champs non-numériques incompatibles et produit les statistiques.
    \item Le Payload JSON sécurisé dresse les tableaux de bord sur l'interface (Recharts).
\end{enumerate}
\lipsum[70-80]

% --- CHAPITRE 4 ---
\chapter{Implémentation et Entraînement des Modèles}
\section{Le Tableau de Bord Interactif}
Intégration de \texttt{recharts} pour tracer les densités et structures via une interface noire, verte émeraude et violette hautement professionnelle.
\lipsum[81-90]

\section{Modèles de Machine Learning}
\subsection{Régression Linéaire Supervisée}
Un pipeline exclusif isole de manière autonome la série visée sur 80\% de données d'entrainements (train set). Les 20\% testés justifient l'indice R2 (\textit{R-Squared}) et l'Erreur Quadratique Moyenne affichées instantanément.

\begin{lstlisting}[language=Python, caption=Exemple d'implémentation du split d'entraînement et de régression extraite du backend]
X_train, X_test, y_train, y_test = train_test_split(X, y, test_size=0.2, random_state=42)
        
model = LinearRegression()
model.fit(X_train, y_train)
predictions = model.predict(X_test)

mse = mean_squared_error(y_test, predictions)
r2 = r2_score(y_test, predictions)
\end{lstlisting}

\lipsum[91-99]

\subsection{K-Means Clustering Non-Supervisé}
Analyse segmentée vectorielle.
\lipsum[100-110]

% --- CHAPITRE 5 ---
\chapter{Bilan, Exportation et Déploiement}
\section{Le Générateur de Rapport PDF Natif}
Plutôt que d'employer une solution CPU-intensive, nous avons codé l'isolation complète du DOM via des directives CSS \texttt{@media print}. Cela rend le navigateur client responsable de tracer et compiler le PDF depuis l'aperçu.
\lipsum[111-120]

\chapter*{Conclusion Générale}
\addcontentsline{toc}{chapter}{Conclusion Générale}
Le produit d'Analyse de Données par IA constitue aujourd'hui un prototype fonctionnel combinant l'agilité des outils Web UI modernes au génie des mathématiques Python. Les enjeux de volumétrie et d'instantanéité ont été gérés, procurant une interface de la plus haute fidélité possible.
\lipsum[121-135]

% --- BIBLIOGRAPHIE ---
\begin{thebibliography}{9}
\bibitem{react}
Documentations Officielles React, META. \url{https://react.dev/}
\bibitem{django}
Django Software Foundation. \url{https://www.djangoproject.com/}
\bibitem{scikit}
Pedregosa et al., Scikit-learn: Machine Learning in Python, JMLR 12, pp. 2825-2830, 2011.
\bibitem{pandas}
Wes McKinney. Data Structures for Statistical Computing in Python, 2010.
\bibitem{tailwind}
Tailwind Labs, A utility-first CSS framework. \url{https://tailwindcss.com/}
\end{thebibliography}

\end{document}
